\chapter{Machbarkeitsanalyse und praktischer Selbstversuch}
\label{ch:explorative_fallstudie}

Nachdem das Experteninterview im vorangegangenen Kapitel die subjektiven Praxisprobleme und die zeitliche Baseline definiert hat, verlagert dieses Kapitel den Fokus auf die technische Ebene der rohen LO-VC-Datenstrukturen. Als zweiter empirischer Baustein zur Analyse des \textit{Environments} nach \textcite{hevnerDesignScienceInformation2004} verfolgt es zwei konkrete Ziele: Zunächst prüft eine Machbarkeitsanalyse (Scoping), welche der idealisierten Praxisanforderungen sich technologisch im Rahmen dieser Arbeit überhaupt automatisieren lassen. Darauf aufbauend dient ein praktischer Selbstversuch dazu, die manuelle Code-Migration Schritt für Schritt zu durchlaufen und exakte syntaktische Transformationsregeln zu extrahieren. Diese praktischen Erkenntnisse liefern die methodische Blaupause, um im Anschluss die Architektur und die System-Prompts des Multi-Agenten-Systems fundiert ableiten zu können.

\section{Technologischer Realitätsabgleich und empirische Scope-Definition}
\label{sec:scope_abgrenzung}

Die im vorherigen Kapitel beschriebenen Praxisanforderungen (vgl. Kapitel~\ref{sec:implikationen}) skizzieren das Idealbild eines zukünftigen Migrationssystems. Nach dem Paradigma des \textit{Design Science Research} \parencite{hevnerDesignScienceInformation2004} dürfen diese Wünsche jedoch nicht ungeprüft in die Systementwicklung übernommen werden. Es ist vorab eine Machbarkeitsprüfung (\textit{Feasibility Check}) notwendig, um den technologischen und funktionalen Umfang (\textit{Scope}) des Prototyps klar abzugrenzen.

Dieser Abschnitt evaluiert die fünf Kernimplikationen aus Kapitel~\ref{sec:implikationen}. Es wird festgelegt, welche Anforderungen im folgenden praktischen Selbstversuch auf Code-Ebene untersucht und anschließend in die Multi-Agenten-Architektur integriert werden:

\subsection{Evaluation der funktionalen Praxisanforderungen}

\begin{enumerate}
\item \textbf{Konsistenzprüfung im Testprozess (Out-of-Scope):}
Die automatisierte Überführung historischer Testkonfigurationen würde Migrationsprojekte erheblich entlasten. Die Umsetzung in diesem Proof-of-Concept ist jedoch \textit{Out-of-Scope}. Eine solche Funktion erfordert eine tiefe Integration in kundenspezifische Testumgebungen und unstrukturiertes Domänenwissen über Altsysteme. Da sich diese Umgebung im Prototyp nicht simulieren lässt, können KI-Agenten die Validierung nicht verlässlich übernehmen. Ein dedizierter „Tester-Agent“ wird daher verworfen.

\item \textbf{Semantische Analyse und Code-Bereinigung (In-Scope):} 
Funktional obsolete Altlasten lassen sich direkt im Code identifizieren. Da diese Bereinigung den Entwickler stark entlastet, wird sie in den Scope aufgenommen. Der Selbstversuch dient dazu, wiederkehrende Muster von „totem Code“ zu erkennen. So können dem zukünftigen Analysten-Agenten entsprechende Erkennungsregeln mitgegeben werden.

\item \textbf{Automatisierte Logik-Transformation (In-Scope -- Kernfokus):} 
Die Baseline zeigt, dass die manuelle Umstrukturierung von Vorbedingungen etwa zwei Drittel des Gesamtaufwands ausmacht. Diese Anforderung bildet den Kern der Systementwicklung. Der Selbstversuch konzentriert sich auf diesen Bereich, um die komplexen Pfade der Code-Umwandlung zu durchdringen und in formale Übersetzungsregeln zu überführen.
\end{enumerate}

\subsection{Evaluation der architektonischen Paradigmen}

\begin{enumerate}
\setcounter{enumii}{3}
\item \textbf{Interaktive Assistenz / Consulting-Ansatz (Fokus auf Explainability):}
Ein frei agierender „Consultant-Agent“, der in einem offenen Dialog strategische Designentscheidungen diskutiert, übersteigt den Rahmen dieses Prototyps. Um dem Praxisbedarf nach Transparenz dennoch gerecht zu werden, wird dieser Ansatz auf das Prinzip der \textit{Explainability} fokussiert. Das System wird so konzipiert, dass es dem Entwickler seine getroffenen architektonischen und syntaktischen Entscheidungen schlüssig in Form von Kommentaren oder kurzen Begründungen darlegt. So bleibt für den Menschen stets nachvollziehbar, warum das System bestimmte Übersetzungspfade gewählt hat. 
\item \textbf{Einhaltung von AVC-Modellierungsrichtlinien (In-Scope):} 
Das System muss standardkonformen AVC-Code erzeugen. Obsolete Elemente, wie klassische Z-Funktionen, müssen zwingend durch native AVC-Standardfunktionen ersetzt werden. Der Selbstversuch hilft dabei, die genauen semantischen Äquivalente im AVC-Standard-Repository zu ermitteln.
\end{enumerate}

\subsection{Synthese des finalen System-Scopes}

Dieser Realitätsabgleich schärft den Fokus der weiteren Forschung. Das Sampling und die qualitative Analyse im folgenden Selbstversuch (vgl. Kapitel~\ref{sec:methodik_fallstudie} ff.) konzentrieren sich folglich ausschließlich auf die \textbf{Code-Bereinigung}, die \textbf{Logik-Transformation} der Vorbedingungen sowie die Integration von \textbf{Explainability} zur Nachvollziehbarkeit der maschinellen Entscheidungen.

Auch die geplante Multi-Agenten-Architektur wird entsprechend eingegrenzt: Das System verzichtet auf autonome Testumgebungen und interaktive Chatbot-Komponenten. Stattdessen fokussiert sich der Prototyp auf ein spezialisiertes, arbeitsteiliges Kernteam von KI-Agenten, welches die Phasen der Datenanalyse, der architektonischen Restrukturierung und der Code-Generierung abdeckt. Die exakte Definition und Ausgestaltung dieser spezifischen Agentenrollen (Version 1) wird direkt aus den Ergebnissen des nun folgenden praktischen Selbstversuchs abgeleitet.

\section{Wissenschaftliche Einordnung und methodisches Vorgehen im Selbstversuch}
\label{sec:methodik_fallstudie}

Das methodische Vorgehen stützt sich auf die Richtlinien für explorative Fallstudien nach \textcite{runesonGuidelinesConductingReporting2009} und die Prinzipien der qualitativen Analyse nach \textcite{seamanQualitativeMethodsEmpirical1999}. In der Softwaretechnik reicht es für die Automatisierung von Prozessen oft nicht aus, nur den ursprünglichen und den finalen Code zu betrachten. Der eigentliche Wert liegt in den kognitiven Zwischenschritten des Entwicklers – dem impliziten Erfahrungswissen (\textit{Tacit Knowledge}), das bei der Problemlösung meist unstrukturiert im Hintergrund abläuft \parencite{seamanQualitativeMethodsEmpirical1999}.

Um dieses Wissen greifbar zu machen, schlüpft der Autor in die Rolle des Entwicklers und migriert reale LO-VC-Daten manuell. Anstatt dabei eine hochkomplexe theoretische Analysemethode anzuwenden, wird der Migrationsprozess systematisch so protokolliert, wie er in der Praxis kognitiv abläuft. Jeder Migrationsschritt wird in drei klar getrennten Phasen dokumentiert:

\begin{enumerate}
\item \textbf{Bestandsaufnahme (Input):} Zunächst wird exakt festgehalten, welche historischen LO-VC-Strukturen und logischen Vorbedingungen für ein Merkmal vorliegen.
\item \textbf{Kognitive Analyse (Gedankengänge):} In diesem Kernschritt wird das \textit{Tacit Knowledge} explizit gemacht. Es wird schriftlich dokumentiert, welche Überlegungen bei der Migration angestellt werden – etwa die Identifikation von veraltetem Code, das Erkennen von Redundanzen oder architektonische Entscheidungen zur Bündelung von Logik.
\item \textbf{Code-Synthese (Output):} Abschließend wird der aus diesen Gedankengängen resultierende, korrekte AVC-Code inklusive eventuell notwendiger Variantentabellen erstellt.
\end{enumerate}

Durch diesen direkten Vergleich von Input, dem eigenen Denkprozess und dem finalen Output lassen sich wiederkehrende Handlungsmuster beim Programmieren identifizieren. Diese Dokumentation der eigenen Denkschritte erfüllt zwei methodische Zwecke: Sie sichert die für Fallstudien geforderte lückenlose Beweiskette (\textit{Chain of Evidence}) \parencite{runesonGuidelinesConductingReporting2009} und bildet die direkte Vorlage, um im nächsten Schritt die konkreten Rollen und Instruktionen (Prompts) für das Multi-Agenten-System abzuleiten.

\section{Datenbasis und Auswahl der Stichprobe}\label{sec:sampling_strategie}

Die Datengrundlage für diesen Selbstversuch bildet ein reales Kundenprojekt der \textit{msg systems ag}. Aufgrund strenger Datenschutzvorgaben, Geschäftsgeheimnissen und stark reglementierter Zugriffsrechte auf historische Altdaten war die Auswahl pragmatisch auf dieses spezifische, verfügbare Projekt beschränkt. Wie \textcite{runesonGuidelinesConductingReporting2009} in ihren methodischen Richtlinien festhalten, ist es in der Praxis der empirischen Softwaretechnik durchaus üblich, Fälle primär nach ihrer reinen Verfügbarkeit auszuwählen.  Aus diesem Projekt wurden fünf Merkmale mitsamt ihrem historisch gewachsenen Beziehungswissen für die manuelle Migration ausgewählt. Die Begrenzung auf genau fünf Merkmale war nicht vorab künstlich definiert, sondern ergab sich aus dem praktischen Verlauf des Selbstversuchs: Nach der Bearbeitung dieser fünf Beispiele zeigte sich, dass sich die grundlegenden kognitiven und syntaktischen Schritte bei der Migration stetig wiederholten. Es traten keine fundamental neuen Muster oder Hürden mehr auf. Mit dem Erreichen dieser inhaltlichen Sättigung war der Zweck des Selbstversuchs erfüllt, da weitere Merkmale keinen zusätzlichen methodischen Mehrwert für die Ableitung der Systemregeln geboten hätten.
\section{Ergebnisse des Selbstversuchs und Ableitung der kognitiven Architektur}
\label{sec:ergebnisse_selbstversuch}

Die vollständigen Protokolle der manuellen Migration – inklusive Input-Code, kognitiven Zwischenschritten und finalem AVC-Code – sind aus Gründen der Lesbarkeit im Anhang \ref{app:migrationsprotokolle} dokumentiert. Dieser Abschnitt fasst die zentralen Erkenntnisse der Fallstudie zusammen und leitet daraus die Handlungsregeln für die Multi-Agenten-Architektur ab.

\subsection{Kognitive Phasen der manuellen Migration}
Die Analyse der Protokolle zeigt, dass die Transformation von LO-VC nach AVC weit mehr ist als ein einfacher Syntax-Austausch. Der Entwickler durchläuft während der Migration unbewusst ein dreistufiges kognitives Modell. Eine bloße 1:1-Übersetzung würde in der neuen AVC-Engine zu fehlerhaftem oder ineffizientem Code führen. Folgende drei Phasen bilden die Blaupause für das zu entwickelnde System:

\begin{enumerate}
\item \textbf{Strukturelle und semantische Bereinigung (Die Analysten-Perspektive)}\
Der historische Kontext muss vor der Transformation bereinigt werden. Die untersuchten Modelle enthielten häufig verwaiste Abhängigkeiten, funktionslosen Legacy-Code, alte Entwicklerkommentare oder logische Redundanzen. Eine rigorose Filterung dieser obsoleten Elemente ist der erste notwendige Schritt, um Fehler im weiteren Verlauf zu vermeiden.

\item \textbf{Architektonische Design-Entscheidungen (Die Architekten-Perspektive)}\\
Nach der Bereinigung muss die strukturelle Zielform festgelegt werden. Dies erwies sich in der Fallstudie als der komplexeste Schritt:
\begin{itemize}
    \item Teilen sich verschiedene Einzelwerte dieselben Vorbedingungen, ist es ineffizient, diese einzeln zu übersetzen. Solche Muster müssen erkannt und zu performanten Variantentabellen aggregiert werden.
    \item Sind Tabellen aufgrund komplexer Ungleichungen nicht optimal, müssen die Werte anders gebündelt werden (z.\,B. durch den \texttt{IN}-Operator).
    \item Zudem muss erkannt werden, welche alten Strukturen (wie der \texttt{INFERENCES}-Block) im AVC gänzlich entfallen.
\end{itemize}

\item \textbf{Syntaktische Synthese (Die Entwickler-Perspektive)}\\
Erst wenn dieser logische Bauplan steht, folgt die eigentliche Programmierung nach den strikten AVC-Regeln. Das umfasst das Schreiben der Restriktionen, die korrekte Pfadreferenzierung (z.\,B. \texttt{\$self.} oder \texttt{\$root.}) und das Ersetzen veralteter Operatoren (z.\,B.\ \texttt{ne} durch \texttt{<>}).
\end{enumerate}

\subsection{Ableitung der V1-Architektur und System-Prompts}

Um diesen komplexen kognitiven Prozess maschinell abzubilden, ist ein einzelnes Sprachmodell (Single-Prompting) ungeeignet. Ein solches Modell wäre mit Bereinigung, Architekturentscheidung und exakter Formatierung gleichzeitig überfordert. Die Ergebnisse des Selbstversuchs erfordern stattdessen eine arbeitsteilige Multi-Agenten-Architektur.

Das in Kapitel 6 konzipierte System wird exakt nach diesem Vorbild aufgebaut: Ein kollaboratives Team aus drei KI-Agenten, unterstützt von einem deterministischen Validierungstool. Aus der Fallstudie lassen sich folgende initiale Rollen und Instruktionen (Prompts) ableiten:

\begin{description}
\item[Rolle 1: LO-VC Data Analyst \& Cleaner] \hfill \
\textbf{Ziel:} Bereinigung der rohen JSON-Eingabedaten.\
\textbf{Handlungsregel:} Analysiere komplexe Vorbedingungen. Identifiziere und ignoriere toten Code und alte Kommentare. Erstelle aus den gültigen Werten eine saubere Liste der steuernden Merkmale.

\item[Rolle 2: AVC Table Architect] \hfill \\
\textbf{Ziel:} Entwurf eines logischen Bauplans.\\
\textbf{Handlungsregel:} Entscheide anhand der bereinigten Daten, welche Merkmale Tabellenspalten bilden. Erkenne identische Vorbedingungen und weise den nachfolgenden Agenten an, diese effizient zusammenzufassen (z.\,B.\ über \texttt{IN()}-Operatoren oder Tabellen).

\item[Rolle 3: AVC Syntax Coder (Synthesizer)] \hfill \\
\textbf{Ziel:} Übersetzung des Bauplans in fehlerfreien AVC-Code.\\
\textbf{Handlungsregel:} Schreibe den finalen Code (\texttt{TABLE}-Aufrufe und \texttt{RESTRICT}-Constraints). Nutze die korrekten Objektreferenzen (wie \texttt{\$self.}) und wende die aktuellen Standard-Operatoren an.

\item[Tool-Integration:] \hfill \\
Übergib den Code abschließend an einen deterministischen Validator. Dieser prüft programmatisch auf verbotene LO-VC-Relikte und erzwingt bei Fehlern eine automatische Korrekturschleife.
\end{description}

Diese empirisch hergeleitete Agentenstruktur bildet das Fundament für die anstehende technische Entwicklung des Systems.